\documentclass[11pt]{article}

\usepackage[margin=0.9in]{geometry}
\usepackage[T1]{fontenc}
\usepackage[utf8]{inputenc}
\usepackage{lmodern}
\usepackage{amsmath,amssymb,amsthm}
\usepackage{booktabs}
\usepackage{enumitem}
\usepackage{longtable}
\usepackage{xcolor}
\usepackage{xurl}
\usepackage[hidelinks]{hyperref}

\setlist[itemize]{leftmargin=1.2em,itemsep=0.25em,topsep=0.25em}
\setlist[enumerate]{leftmargin=1.4em,itemsep=0.25em,topsep=0.25em}
\setlength{\parindent}{0pt}
\setlength{\parskip}{0.55em}
\emergencystretch=2em
\sloppy

\newcommand{\lean}[1]{\path{#1}}
\newcommand{\R}{\mathbb{R}}
\newcommand{\E}{\mathbb{E}}
\newcommand{\Prob}{\mathbb{P}}
\newcommand{\toP}{\xrightarrow{p}}
\newcommand{\toD}{\xrightarrow{d}}
\newcommand{\op}{o_p}
\newcommand{\Op}{O_p}
\newcommand{\iid}{\mathrm{i.i.d.}}
\newcommand{\Var}{\operatorname{Var}}
\newcommand{\Cov}{\operatorname{Cov}}
\newcommand{\rank}{\operatorname{rank}}

\title{Issue 42 Proof Summary: Consolidated Chapter 2--7 Theorem Progress}
\author{Hansen Econometrics Lean Formalization}
\date{Branch \lean{codex/issue-42-ch7-unblocked}, compiled May 9, 2026}

\begin{document}
\maketitle

\begin{abstract}
This document is a self-contained mathematical map of the results added in the
consolidated issue-42 branch.  The branch is large because it collects many
small theorem slices that were first developed as separate checkpoints.  This
summary explains the mathematical surface now present in Lean, records the main
public theorem names, and separates completed theorem coverage from remaining
gaps.  It is not a replacement for the Lean files; it is intended to make the
pull request reviewable.
\end{abstract}

\tableofcontents

\section{Conventions and Proof Architecture}

The formalization keeps three recurring conventions.

\begin{itemize}
  \item Convergence in probability is written in Lean with
  \lean{TendstoInMeasure}; convergence in distribution is written with
  \lean{TendstoInDistribution}.
  \item Finite-sample ordinary least squares has a base API requiring an
  invertibility typeclass.  Large-sample theorem statements use totalized
  estimators such as \lean{olsBetaStar} and textbook-facing wrappers such as
  \lean{olsBetaOrZero}.  The star estimator uses the nonsingular inverse and is
  total; the OrZero wrapper is the reader-facing branch that agrees with
  ordinary OLS when the Gram matrix is nonsingular.
  \item Many Chapter 7 theorems are stated at a condition-package layer.  The
  branch adds direct compact and iid sufficient-condition packages for several
  robust-HC endpoints, but some assumption surfaces remain stronger than the
  literal textbook assumptions.
\end{itemize}

The key organizing idea is to prove reusable asymptotic infrastructure once in
Chapter 6 and then reuse it in Chapter 7.  The Chapter 7 results are mostly
composition theorems: consistency plus score CLTs, continuous mapping, Slutsky,
matrix inverse continuity, and Gaussian linear-image laws.

\section{Chapter 2: Potential Outcomes and CIA}

The branch adds a variable-facing potential-outcomes layer for Hansen's causal
notation.  Given a treatment indicator $D$, untreated and treated potential
outcomes $Y(0),Y(1)$, and covariates $X$, the Lean API defines:
\[
  Y = D Y(1) + (1-D)Y(0), \qquad
  ACE = \E[Y(1)-Y(0)],
\]
and a variable-conditioned CATE
\[
  ACE(X) = \E[Y(1)-Y(0)\mid X].
\]

Main declarations:
\begin{itemize}
  \item \lean{observedOutcome}, \lean{treatmentEffect},
  \lean{averageTreatmentEffect}, and
  \lean{conditionalAverageTreatmentEffectOn}.
  \item \lean{PotentialOutcomeCIAOn}: the conditional-independence assumption
  for potential outcomes, treatment, and covariates.
  \item \lean{PotentialOutcomeCIAOn.toTreatmentMeanIndependentOn}: CIA implies
  the mean-independence condition needed for regression-style CATE statements.
  \item \lean{condExpOn_observedOutcome_treatment_covariates_eq_branch_of_CIA}:
  the observed-outcome regression on $(D,X)$ selects the treated or untreated
  branch under CIA.
  \item \lean{conditionalAverageTreatmentEffectOn_treatment_covariates_eq_of_CIA}:
  conditioning the treatment effect on $(D,X)$ gives the same CATE as
  conditioning on $X$ alone, almost everywhere.
  \item \lean{conditionalAverageTreatmentEffectSurface},
  \lean{observedRegressionSurface}, and
  \lean{observedRegressionTreatmentContrastSurface}: thin pointwise notation
  bridges for $ACE(x)$, $m(d,x)$, and $m(1,x)-m(0,x)$.
  \item \lean{conditionalAverageTreatmentEffectOn_eq_surface} and
  \lean{conditionalAverageTreatmentEffectOn_eq_observedRegressionTreatmentContrastSurface}:
  supplied pointwise versions pull back to the existing variable-facing CATE.
\end{itemize}

The result is a usable a.e. version of the potential-outcomes and conditional
average-treatment-effect story with pointwise notation bridges layered on top
of the a.e. conditional-expectation core.

\section{Chapter 3: Finite-Sample Least-Squares Algebra}

Chapter 3 receives finite-sample algebra useful for later asymptotics and
diagnostics.

\subsection{Partitioned Regression and FWL}

For a column partition $X=[X_1\;X_2]$, the branch formalizes the Frisch--Waugh--
Lovell coefficient formulas.  The central public theorem is
\lean{fromColsBeta_eq_partitionedBetaFormulas}, which packages the two block
coefficient identities.  The supporting definitions
\lean{partitionedLeftBetaFormula} and \lean{partitionedRightBetaFormula}
match Hansen's displayed formulas for the two coefficient blocks.

\subsection{Projection, Leverage, and Influence}

The branch adds leverage and leave-one-out identities.  The formal statements
include:
\begin{itemize}
  \item \lean{leaveOneOutBetaDeleted_eq_leaveOneOutBeta}: the literal row-deleted
  regression agrees with the reduced Gram/cross-product formula.
  \item \lean{leaveOneOutBeta_eq_olsBeta_sub_invGram_mulVec}: the leave-one-out
  coefficient differs from the full-sample coefficient by a leverage-weighted
  correction.
  \item \lean{leaveOneOutResidual_eq_inv_one_sub_leverage_mul_residual}: the
  leave-one-out residual is $e_i/(1-h_{ii})$.
  \item \lean{fitted_sub_leaveOneOutPrediction_eq_leverage_mul_residual}: the
  prediction change equals the leverage times the leave-one-out residual.
  \item \lean{maxPredictionInfluence_eq_maxLeveragePredictionErrorInfluence}:
  the maximum leave-one-out prediction influence has the leverage-residual
  expression advertised by the diagnostics.
\end{itemize}

These are deterministic matrix identities.  They later feed Chapter 7 leverage
and residual-rate statements.

\section{Chapter 4: Finite-Sample Covariance Estimators}

The branch extends the finite-sample covariance-estimator layer that Chapter 7
uses asymptotically.

\subsection{HC Ordering}

For leverage values $0\le h_{ii}<1$, the branch proves the positive-semidefinite
ordering
\[
  HC0 \preceq HC2 \preceq HC3.
\]
The public matrix-level results are
\lean{olsHuberWhiteVarianceEstimator_le_HC2_posSemidef} and
\lean{olsHuberWhiteHC2VarianceEstimator_le_HC3_posSemidef}.

\subsection{HC2 and HC3 Conditional Expectations}

The branch records exact conditional-expectation formulas for leverage-adjusted
residual squares:
\begin{itemize}
  \item \lean{condExp_HC2_adjusted_annihilator_row_sq_eq_sigmaSq} and
  \lean{condExp_HC2_adjusted_annihilator_row_sq_eq_diagonal}.
  \item \lean{condExp_HC3_adjusted_annihilator_row_sq_eq_sigmaSq_mul_inv} and
  \lean{condExp_HC3_adjusted_annihilator_row_sq_eq_diagonal}.
  \item Matrix-valued expectation wrappers for the corresponding HC2 and HC3
  covariance estimators.
\end{itemize}

The branch also adds singleton-cluster bridges showing that singleton clustered
covariance estimators reduce to the heteroskedastic sandwich estimators, and
tightens the finite-sample Gauss--Markov proof layer.

\section{Chapter 6: Asymptotic Infrastructure}

Chapter 6 is the main infrastructure layer for issue 42.  The branch adds or
extends large-sample theorem wrappers that Chapter 7 reuses.

\subsection{Convergence in Probability and Distribution}

The core CMT and WLLN declarations are:
\begin{itemize}
  \item \lean{tendstoInMeasure_continuous_comp}: global CMT in probability.
  \item \lean{tendstoInMeasure_continuousAt_const_comp}: CMT at a point for a
  constant limit.
  \item \lean{tendstoInDistribution_continuous_comp}: global CMT in
  distribution.
  \item \lean{probabilityMeasure_tendsto_map_of_tendsto_of_ae_continuous}
  and \lean{tendstoInDistribution_ae_continuous_comp}: the textbook
  a.s.-continuity CMT in distribution, proved by a Portmanteau closed-set
  push-forward argument.
  \item \lean{tendstoInMeasure_wlln}: Banach-valued iid WLLN, obtained by
  upgrading Mathlib's strong law to convergence in measure.
  \item \lean{tendstoInMeasure_transformed_wlln}: transformed WLLN for
  sample means of $h(Y_i)$.
\end{itemize}

The branch also adds product, sum, matrix-vector, matrix-multiplication, matrix
inverse, coordinatewise-product, and finite-sum convergence lemmas used in
Chapter 7.

\subsection{Cramer--Wold and CLTs}

The branch formalizes the Cramer--Wold bridge and iid finite-dimensional CLTs:
\begin{itemize}
  \item \lean{cramerWold_tendstoInDistribution}: convergence of all scalar
  projections implies vector convergence in distribution.
  \item \lean{iidScalarCLT_tendstoInDistribution_gaussian}: scalar iid CLT
  wrapper over Mathlib.
  \item \lean{iidVectorProjectionCLT_tendstoInDistribution_gaussian_cov}:
  scalar projection CLT with covariance identified by the vector covariance
  matrix.
  \item \lean{iidVectorCLT_tendstoInDistribution_multivariateGaussian}:
  multivariate iid CLT with limit $N(0,\operatorname{covMat}(Y_0))$.
\end{itemize}

The branch also introduces projection-level endpoints for heterogeneous and
Lindeberg triangular-array CLTs:
\lean{MultivariateLindebergCLTConditions},
\lean{multivariateLindebergCLT_tendstoInDistribution},
\lean{HeterogeneousArrayCLTConditions}, and
\lean{heterogeneousArrayCLT_tendstoInDistribution}.  These are Cramer--Wold
endpoints; the scalar triangular-array probability engine remains outside this
branch.

\subsection{Delta Method}

The new Delta-method module proves the deterministic remainder statement
\[
  g(x)-g(\theta)-Dg(\theta)(x-\theta) = o(\|x-\theta\|)
\]
from Frechet differentiability and packages the distributional conclusion:
\[
  p_n(\hat\theta_n-\theta)\toD Z
  \quad\Longrightarrow\quad
  p_n(g(\hat\theta_n)-g(\theta))\toD Dg(\theta)Z,
\]
assuming the linearization remainder is $o_p(1)$.

Main declarations:
\begin{itemize}
  \item \lean{deltaMethod_remainder_isLittleO}.
  \item \lean{deltaMethod_tendstoInDistribution}.
  \item \lean{smoothFunction_consistency}.
  \item \lean{smoothFunction_asymptoticNormality_gaussian}.
\end{itemize}

Concrete examples cover coordinate square, inverse, product, and ratio maps,
with derivative matrices and Gaussian linear-image laws:
\lean{coordinateSquareDerivativeMatrix_hasLaw_multivariateGaussian_zero},
\lean{coordinateInvDerivativeMatrix_hasLaw_multivariateGaussian_zero},
\lean{coordinateProductDerivativeMatrix_hasLaw_multivariateGaussian_zero}, and
\lean{coordinateRatioDerivativeMatrix_hasLaw_multivariateGaussian_zero}.

\subsection{Stochastic Order}

The branch adds a local stochastic-order layer:
\begin{itemize}
  \item \lean{BoundedInProbability}: scalar $O_p(1)$.
  \item \lean{BoundedInProbabilityNorm}: norm-valued $O_p(1)$.
  \item moment-to-$O_p$ and moment-to-$o_p$ bridges, including scaled first
  moments, natural powers, and positive real powers.
  \item closure under addition, multiplication, deterministic scaling, and
  ``small times bounded is small''.
\end{itemize}

These statements are the workhorse behind Chapter 7 remainder terms, inverse
gaps, max-row rates, residual rates, and leverage rates.

\subsection{Uniform Integrability, Maximum Bounds, and Weak Moments}

The maximum theorem for row norms is formalized on the nonnegative power scale:
\[
  n^{-1}\max_{i<n} Z_i \toP 0
\]
from uniform integrability of $Z_i$.  In applications $Z_i=\|X_i\|^r$, and a
root CMT gives Hansen's displayed maximum-rate form.  Main declarations are
\lean{uniformIntegrable_tail_eLpNorm_one},
\lean{uniformIntegrable_one_iff_tail_eLpNorm},
\lean{uniformIntegrable_one_of_identDistrib_memLp},
\lean{max_norm_scaled_tendstoInMeasure_zero_of_uniformIntegrable_norm_r}, and
\lean{max_norm_scaled_tendstoInMeasure_zero_of_identDistrib_memLp}.

For weak convergence and moments, the branch adds:
\begin{itemize}
  \item bounded-continuous Portmanteau moment wrappers;
  \item norm-truncation functions and weak-convergence norm-bound transfer;
  \item limit integrability under eventual first-moment bounds;
  \item real clipping functions and clipped-identity weak-moment convergence;
  \item pointwise and integral bounds showing clipping error is controlled by
  large-tail integrals;
  \item conversion from Mathlib's \lean{UniformIntegrable} tail seminorms to
  real clipping-tail bounds, and the resulting weak-convergence/UI moment
  theorem.
\end{itemize}

The main theorem names are
\lean{TendstoInDistribution.integral_norm_limit_le_of_eventually_integral_norm_bound},
\lean{realClipBoundedContinuousFunction},
\lean{TendstoInDistribution.integral_realClip_tendsto},
\lean{abs_integral_sub_realClip_le_two_mul_integral_tail_abs}, and
\lean{TendstoInDistribution.integral_tendsto_of_uniformIntegrable}.

\subsection{Sample-Mean Sharpness}

The branch adds exact scalar and vector sample-mean variance formulas:
\begin{itemize}
  \item \lean{iidSampleMean_variance_eq_inv_card_mul}.
  \item \lean{iidLinearUnbiasedEstimator_variance_ge_sampleMean}.
  \item \lean{iidLinearUnbiasedEstimator_covMat_sub_sampleMean_posSemidef}.
  \item \lean{iidSampleMean_covMat_eq_inv_card_smul}.
  \item \lean{iidSampleMean_covMat_eq_inv_card_smul_of_identDistrib}.
\end{itemize}

These prove the sample mean attains the $n^{-1}$ variance/covariance scale for
iid observations and that no scalar or finite-dimensional scalar-weighted linear
unbiased estimator improves on the equal-weight sample mean.  They are not the full
arbitrary-unbiased-estimator lower bound of Hansen Theorem 6.11.

\section{Chapter 7: Least-Squares Asymptotics}

The Chapter 7 additions assemble the Chapter 6 infrastructure into OLS
asymptotic statements.

\subsection{Condition Packages}

The public condition structures now include:
\begin{itemize}
  \item \lean{LeastSquaresConsistencyConditions};
  \item \lean{ErrorVarianceConsistencyConditions};
  \item \lean{ScoreCLTConditions};
  \item \lean{RobustCovarianceConsistencyConditions};
  \item \lean{HomoskedasticErrorVariance};
  \item \lean{RobustFeasibleHCMomentConditions};
  \item \lean{IidRobustFeasibleHCMomentConditions}.
\end{itemize}

The branch adds conversions among these packages and direct compact/iid
sufficient-condition endpoints for many theorem statements.  These packages are
review aids as well as proof infrastructure: they isolate exactly what is being
assumed for each theorem surface.

\subsection{OLS Consistency and Normality}

The branch proves consistency and asymptotic normality for the totalized and
ordinary-wrapper estimators.  Key statements include:
\begin{itemize}
  \item \lean{olsBetaStar_consistent} and OrZero/literal wrappers.
  \item \lean{scoreVector_sampleCrossMoment_tendstoInDistribution_multivariateGaussian}.
  \item \lean{feasibleScore_tendstoInDistribution_scoreCLT}.
  \item \lean{olsBetaStar_vector_tendstoInDistribution_scoreCLT}.
  \item \lean{olsBetaStar_vector_tendstoInDistribution_multivariateGaussian}.
  \item \lean{olsBetaOrZero_vector_tendstoInDistribution_multivariateGaussian}.
  \item \lean{scoreProj_olsBetaStar_tendstoInDistribution_gaussian_cov_all}.
\end{itemize}

Mathematically, these statements show
\[
  \sqrt{n}(\hat\beta_n-\beta)\toD N(0,Q^{-1}\Omega Q^{-1})
\]
at the current condition layer, with projection-family and vector forms.

\subsection{Residual Variance and Homoskedastic Covariance}

For Hansen Theorems 7.4 and 7.5, the branch proves consistency of totalized
residual-variance and homoskedastic plug-in covariance estimators:
\begin{itemize}
  \item \lean{olsSigmaSqHatStar_tendstoInMeasure_errorVariance}.
  \item \lean{olsS2Star_tendstoInMeasure_errorVariance}.
  \item \lean{olsHomoCovStar_tendstoInMeasure}.
  \item direct iid endpoints under
  \lean{IidRobustFeasibleHCMomentConditions}.
\end{itemize}

\subsection{HC0, HC1, HC2, and HC3 Covariance Consistency}

For Theorems 7.6 and 7.7, the branch adds a substantial robust-HC consistency
layer:
\begin{itemize}
  \item HC0 middle-matrix and sandwich consistency for feasible residuals.
  \item HC1 consistency using the finite-sample scale factor
  \lean{hc1FiniteSampleScale_tendsto_one}.
  \item HC2 and HC3 consistency through a generic leverage-adjusted middle
  matrix and max-leverage control.
  \item compact moment and iid joint-observation endpoints for HC0--HC3.
\end{itemize}

Representative theorem names are
\lean{olsHetCovStar_tendstoInMeasure_of_robustFeasibleHCMomentConditions},
\lean{olsHetCovHC1Star_tendstoInMeasure_of_robustFeasibleHCMomentConditions},
\lean{olsHetCovHC2Star_tendstoInMeasure_of_robustFeasibleHCMomentConditions},
\lean{olsHetCovHC3Star_tendstoInMeasure_of_robustFeasibleHCMomentConditions},
and the corresponding
\lean{..._of_iidRobustFeasibleHCMomentConditions} endpoints.

\subsection{Functions of Parameters}

For Theorems 7.8--7.10, the branch proves both fixed-linear and nonlinear
function-of-parameter results:
\begin{itemize}
  \item continuous-map consistency for $r(\hat\beta)$;
  \item nonlinear Delta-method Gaussian limits for totalized and OrZero OLS;
  \item fixed-linear and random-linear covariance CMTs;
  \item plug-in derivative covariance consistency for nonlinear functions.
\end{itemize}

Representative names include
\lean{nonlinearFunction_olsBetaStar_delta_tendstoInDistribution_gaussian},
\lean{nonlinearFunction_olsBetaOrZero_delta_tendstoInDistribution_gaussian},
\lean{linMapCov_tendstoInMeasure},
\lean{randomLinearMapCovariance_tendstoInMeasure},
\lean{nonlinearDerivativeCovariance_olsBetaStar_tendstoInMeasure}, and
\lean{nonlinearDerivativeCovariance_olsBetaOrZero_tendstoInMeasure}.

\subsection{Standard Errors, t Statistics, Confidence Intervals, and Wald Tests}

For Theorems 7.11--7.14, the branch adds the inference layer:
\begin{itemize}
  \item standard-error consistency by continuous mapping;
  \item studentized scalar limits;
  \item symmetric confidence-interval coverage from absolute t-statistic
  convergence;
  \item multivariate Wald statistic convergence;
  \item homoskedastic and heteroskedastic HC0--HC3 versions;
  \item direct compact and iid endpoints where the required condition packages
  can be discharged.
\end{itemize}

Core declarations include
\lean{studentizedLimit_tendstoInDistribution},
\lean{linMapCovStdError_tendstoInMeasure},
\lean{symmetricCI_coverage_of_abs_tstat_nonpos_tendsto_zero},
\lean{symmetricCI_coverage_of_abs_tstat_standardNormal_se_tendsto_pos},
\lean{waldQuadForm_tendstoInDistribution_of_vector_and_covariance},
\lean{hasLaw_gaussian_mahalanobis_chiSquared}, and the HC0--HC3 scalar and
multivariate wrappers in \lean{Inference.lean}, \lean{Normality.lean}, and
\lean{SandwichAssembly.lean}.

\subsection{Residual Uniformity and Leverage Rates}

For Theorems 7.16 and 7.17, the branch proves deterministic inequalities and
probabilistic rate wrappers:
\begin{itemize}
  \item the max residual error is bounded by a coefficient-error factor times a
  max-row-norm factor;
  \item the coefficient-error factor is $O_p(1)$ by the vector OLS CLT;
  \item the max-row-norm factor is $o_p(1)$ by the Chapter 6 maximum theorem;
  \item max leverage is controlled by inverse-Gram bounds and max row norms.
\end{itemize}

Representative theorem names are
\lean{sqrt_smul_olsBetaStar_sub_boundedInProbabilityNorm},
\lean{sqrt_scaledMaxRowNorm_sq_tendstoInMeasure_zero_of_uniformIntegrable_norm_sq},
\lean{maxResidualErrorStar_tendstoInMeasure_zero_of_uniformIntegrable_rowNorm_sq},
\lean{maxResidualErrorStar_tendstoInMeasure_zero_of_identDistrib_memLp_rowNorm_sq},
\lean{scaledMaxLeverageStar_tendstoInMeasure_zero_of_scaled_maxRowNorm_sq},
\lean{maxLeverageStar_tendstoInMeasure_zero_of_uniformIntegrable_rowNorm_sq}, and
\lean{maxLeverageStar_tendstoInMeasure_zero_of_identDistrib_memLp_rowNorm_sq}.
The direct iid feasible-HC package derives the squared-row $L^1$ input for
these wrappers from its fourth-row-moment assumption. It also derives the
regressor outer-product $L^1$ field from the same fourth-row moment, and derives
the score-vector $L^1$ field and the score-projection $L^2$ field used by the
CLT package from the score outer-product moment. The combined robust feasible-HC
package and the iid package now also derive the feasible-HC third absolute row-weight moment
$\mathbb{E}\{|e_0|\|X_0\|^3\}<\infty$ from the score outer-product moment and
the fourth-row moment, rather than carrying it as a separate primitive field in
those public packages.

\subsection{Edgeworth Interface}

The branch does not prove Hansen's concrete high-moment Edgeworth expansion.
It does add generic interfaces that future work can instantiate:
\begin{itemize}
  \item \lean{FirstOrderEdgeworthExpansion};
  \item \lean{SecondOrderEdgeworthExpansion};
  \item scaled-remainder consequences;
  \item unscaled second-order approximation-error consequences;
  \item CDF convergence consequences;
  \item a bridge from the second-order interface to the first-order interface.
\end{itemize}

Thus Theorem 7.15 has a reusable formal target shape, but not the cumulant or
polynomial expansion proof.

\section{Remaining Gaps}

The branch is not a claim that Chapters 2--7 are literally complete.  The main
remaining gaps are:

\begin{longtable}{p{0.16\textwidth}p{0.77\textwidth}}
\toprule
Area & Remaining work \\
\midrule
Chapter 4 & Clustered covariance asymptotics and multi-observation cluster
partition infrastructure; HC4-style leverage-adjusted variants. \\
Chapter 5 & Kinal theorem 5.5, deferred unless a downstream theorem needs the
exact moment-existence threshold. \\
Chapter 6 & Scalar triangular-array Lindeberg/Lyapunov sufficient-condition
discharge for Theorems 6.4 and 6.5; arbitrary-unbiased-estimator lower bound
for Theorem 6.11. \\
Chapter 7 & Further tightening of the remaining public condition packages to the
literal textbook iid/moment assumptions; concrete cumulant/polynomial Edgeworth
expansion for Theorem 7.15. \\
\bottomrule
\end{longtable}

\section{Files Touched by the Consolidated Branch}

The main proof files changed by the branch are:
\begin{itemize}
  \item \lean{HansenEconometrics/AsymptoticUtils.lean}
  \item \lean{HansenEconometrics/AsymptoticUtils/DeltaMethod.lean}
  \item \lean{HansenEconometrics/AsymptoticUtils/MaxBounds.lean}
  \item \lean{HansenEconometrics/AsymptoticUtils/StochasticOrder.lean}
  \item \lean{HansenEconometrics/Chapter2PotentialOutcomes.lean}
  \item \lean{HansenEconometrics/Chapter3FWL.lean}
  \item \lean{HansenEconometrics/Chapter3LeastSquaresAlgebra.lean}
  \item \lean{HansenEconometrics/Chapter3Projections.lean}
  \item \lean{HansenEconometrics/Chapter4LeastSquaresRegression.lean}
  \item \lean{HansenEconometrics/Chapter6Asymptotics.lean}
  \item \lean{HansenEconometrics/Chapter7Asymptotics/Consistency.lean}
  \item \lean{HansenEconometrics/Chapter7Asymptotics/Inference.lean}
  \item \lean{HansenEconometrics/Chapter7Asymptotics/MiddleConsistency.lean}
  \item \lean{HansenEconometrics/Chapter7Asymptotics/Normality.lean}
  \item \lean{HansenEconometrics/Chapter7Asymptotics/SampleMiddle.lean}
  \item \lean{HansenEconometrics/Chapter7Asymptotics/SandwichAssembly.lean}
  \item \lean{HansenEconometrics/ProbabilityUtils.lean}
\end{itemize}

The chapter inventories were updated alongside the Lean files, and this document
is intended as a single reviewable guide to the mathematical content.

\end{document}
